\begin{afterword}
\summaryhead

The post argues that evolution does not work for the good of the species. Selection changes gene frequencies, and a gene spreads according to what it does for its own copies. The sex ratio is the main example. A group with more daughters would grow faster, but when males are rare each son is worth more to the parent, so individual selection keeps investment in sons and daughters equal. Near-equal human births show that individual selection wins.

From this the post argues that a species can evolve to extinction. A gene that sacrifices its carrier to save the species loses frequency. Viruses may kill their hosts too fast, a sex-ratio distorter in mice could drive them extinct, transposons fill genomes, a perfect DNA-copying mechanism would spread and leave the species unable to adapt, and the bystander effect may be the product of an evolutionary arms race. Cancer is evolution among cells that kills the body. The author suggests humanity may be evolving to extinction now.

\responsehead

The post's main claim is standard biology, and it explains it well. Selection does not act for the good of the species, and it can favour genes that harm or doom the population that carries them. The ideas are older than the post. The critique of reasoning about the good of the species is George C.~Williams's (1966), the sex-ratio argument is usually called Fisher's principle, and the result for a driving Y chromosome is Hamilton's (1967). Later literature calls the phenomenon evolutionary suicide. The post names none of these, which a short essay need not do, but a reader should know where the arguments come from.

The weakness is the evidence that species actually evolve to extinction. The post's cases fall into three groups. The Frodo gene and the perfect copier are thought experiments. The virus case (``probably happened any number of times''), the mice (``if mice were to evolve the ability to fly'') and humanity (``may be finishing up'') are hedged guesses, and the transposons ``may not extinguish a species.'' Only cancer is an observed case, and it concerns cells in one body. The hedges are honest, and the thought experiments do show that extinction by selection is possible. But the post gives no observed case for a species or population, although one existed: in a 2003 laboratory study, cheating bacteria drove whole populations extinct.

The mouse example could not be traced. The post describes a mouse gene on the male sex chromosome that yields only sons, and no such gene was found. The mouse case to which group selection was applied is the t haplotype, on chromosome 17, which carrier males pass to most of their offspring, and which in two copies causes male sterility or death of the embryo. The post's reasoning fits a driving Y chromosome, which Hamilton analysed; the facts it attaches to mice do not.

The ending moves furthest from its evidence. A laboratory study of students reporting a seizure over an intercom becomes a speculated genetic arms race, and then a claim that few people are ``steppin' forward'' against humanity's extinction threats. The speculation is flagged as such. The claim about humanity has no support in the post.

In short: a correct account, on the lines of Williams, Fisher and Hamilton, of how selection can harm a species, whose cases of actual extinction are hedged guesses, a mouse gene that could not be traced, and a speculation about humanity.
\end{afterword}
